% Generated draft: requires mathematical and source review.
\documentclass[11pt]{article}
\usepackage[margin=25mm]{geometry}
\usepackage{fontspec}
\setmainfont{Latin Modern Roman}
\usepackage{amsmath,amssymb,mathtools,unicode-math}
\setmathfont{Latin Modern Math}
\usepackage{enumitem,needspace,xurl,hyperref}
\setlength{\parindent}{0pt}
\setlength{\parskip}{5pt}
\setlength{\emergencystretch}{3em}
\newcommand{\N}{\mathbb{N}}
\newcommand{\Z}{\mathbb{Z}}
\newcommand{\Q}{\mathbb{Q}}
\newcommand{\R}{\mathbb{R}}
\newcommand{\C}{\mathbb{C}}
\newcommand{\F}{\mathbb{F}}
\newcommand{\A}{\mathbb{A}}
\newcommand{\PP}{\mathbb{P}}
\newcommand{\E}{\mathbb{E}}
\newcommand{\eps}{\varepsilon}
\newcommand{\dd}{\,\mathrm{d}}
\begin{document}
\section*{UnsolvedMath ID: 701714}
\textbf{Source problem code:} catalogue entry
\subsection{Definitions and mathematical statement}
A group \(G\) is periodic if every element of \(G\) has finite order. It is locally finite if every subgroup generated by finitely many elements is finite. Let
\[
I(G)=\{t\in G:t^2=1,\ t\ne 1\}
\]
be the set of involutions of \(G\). For \(t\in G\), define its centralizer by
\[
C_G(t)=\{g\in G:gt=tg\}.
\]
The condition that \(G\) is generated by involutions means \(G=\langle I(G)\rangle\); no finiteness assumption is imposed on this generating set.

Is it true that for every periodic group \(G\),
\[
\left(G=\langle I(G)\rangle\ \text{and}\ \forall t\in I(G),\ C_G(t)\text{ is locally finite}\right)
\ \Longrightarrow\ G\text{ is locally finite}?
\]
The supplied registry labels this question open; that status is not independently verified here.
\subsection{Short English statement}
Must every periodic group generated by involutions be locally finite if the centralizer of each involution is locally finite?
\subsection{Sources}
\begin{itemize}
\item Local UnsolvedMath registry, supplied complete source record: permanent ID \(701714\), exact source problem code \texttt{catalogue entry}, title \emph{Kourovka Notebook Problem 20.56}.
\item The supplied record attributes Problem 20.56 to V. D. Mazurov, proposed in 2022, and cites \emph{The Kourovka Notebook (Unsolved Problems in Group Theory)}, 20th issue (2022), and the 21st issue edition (2026), identified as arXiv:1401.0300v48. These bibliographic details are supplied attribution and have not been independently verified.
\end{itemize}
\end{document}
